\section*{Experience}

\cventry
    {Argonne National Laboratory}
    {Lemont, IL}
    {Research Aide Technical -- PhD}
    {May--Aug.\ 2026; May--Aug.\ 2025}

\begin{itemize}
    \item Co-developed VisKnacks, an open-source system that turns
        natural-language requests into ParaView visualization workflows.
        Integrated agent orchestration, documentation retrieval, and
        containerized rendering so researchers could refine workflows
        locally and render large simulation datasets on remote HPC systems.
    \item Benchmarked three large language models on 10 ParaView tasks
        on Argonne's Crux system, comparing rendered outputs with
        human-generated reference images using PSNR and SSIM.
        First-authored and presented the VisKnacks poster at SC26.
    \item Co-developed SCODA, a reproducible Docker Compose benchmarking
        pipeline for supercomputing operations analytics. Worked with
        Argonne's business intelligence team to turn operational workflows
        into ingestion, concurrent read/write, and aggregation benchmarks
        using real and synthetic system logs.
    \item Evaluated nine distributed database and lakehouse platforms
        using throughput, latency, and resource utilization to identify
        workload-specific tradeoffs. First-authored and presented the
        SCODA poster at SC25.
\end{itemize}

\cventry
    {Loyola University Chicago}
    {Chicago, IL}
    {Teaching Assistant, Department of Computer Science}
    {Aug.\ 2024--Present}

\begin{itemize}
    \item Lead software engineering activities for the Software and
        Systems Laboratory, coordinating across disciplines to resolve
        issues spanning research and software development.
    \item Hold office hours, grade assignments and exams, and proctor
        exams across five computer science courses, including Machine
        Learning, Software Engineering, and Computer Security.
\end{itemize}

\cventry
    {Loyola University Chicago}
    {Chicago, IL}
    {Graduate Research Assistant, Department of Computer Science}
    {Aug.\ 2022--May 2023; Aug.\ 2023--Aug.\ 2024}

\begin{itemize}
    \item Investigated how developers discover and reuse pretrained
        models on Hugging Face and how model dependencies enter open-source
        software supply chains. Co-authored an empirical study on model
        reuse (ICSE 2023) and the PeaTMOSS dataset paper (MSR 2024).
    \item Mentored graduate researchers on applying AI and machine
        learning to software engineering research and documented
        workflows and troubleshooting steps so peers could reproduce
        and extend the work.
\end{itemize}

\cventry
    {Argonne National Laboratory}
    {Lemont, IL}
    {Research Aide Technical -- Master's}
    {May--Aug.\ 2023}

\begin{itemize}
    \item Used Argonne supercomputing resources to trace how pretrained
        deep learning models were developed and reused downstream.
        Collaborated on methods to process model artifacts and compare
        pretrained models empirically.
\end{itemize}

\cventry
    {Loyola University Chicago}
    {Chicago, IL}
    {Undergraduate Research Assistant, Department of Computer Science}
    {Sept.\ 2020--May 2022}

\begin{itemize}
    \item Co-led a four-person team that built SSL Metrics, an open-source
        Python toolkit that uses Git and GitHub data to track defect
        density, issue age, and development activity over time.
    \item Configured Travis CI to convert reStructuredText content and
        deploy it to the lab's production website.
\end{itemize}

\cventry
    {Medline Industries}
    {Mundelein, IL}
    {IS/IT Developer Analyst}
    {May 2021--May 2022}

\begin{itemize}
    \item Built internal talent identification and workforce planning
        applications with OutSystems and Oracle Database, translating
        requirements from IT and business stakeholders into features
        through Agile development cycles.
    \item Troubleshot production application issues across projects and
        business groups as part of a 24/7 on-call rotation.
\end{itemize}